% Contribution /root/noncm_gate; model GPT-6 Astra Pro; effort not exposed.
% Whole-proof check events are supplied separately in CHECK_EVENTS.tsv.
% Bibliography entries are supplied in bibliography.tex.

\section{The non-CM gate for curves and abelian varieties}
\label{sec:noncm-gate-next}

We use the convention of ZETA\_TWO CM01:
$\operatorname{wt} M(n)=\operatorname{wt}M-2n$ and
$L(M(n),s)=L(M,s+n)$.  ``CM'' in this section means
\emph{potentially CM in the category of motives for absolute Hodge
cycles}, including every Artin motive.  In particular a nonabelian
finite-image Artin representation is in the CM category.  The splitting
results below are unconditional and do not use either zero theorem.
Only the assertions about the zeros of their CM factors use the
conditional input \textbf{M} of ZETA\_TWO CM01.  Its scope is all unitary
Hecke characters, including infinite-order characters; the source and
dependency audit of that input are carried from CM01 and MDOC01.
For a chosen finite-type integral model $\mathcal X$, the superscript
$S$ means the Euler product over $p\notin S$.  The full model zeta
function is always
\[
 \zeta_{\mathcal X}(s)=
       \left(\prod_{p\in S}Z(\mathcal X_p,p^{-s})\right)
                                  \zeta_X^S(s).
\]
The actual exceptional fibres therefore remain explicit; none is
replaced by a generic-fibre inertia factor without the correction in
ZETA\_TWO ZT04.

\subsection{NG01: a canonical maximal-CM summand, with descent}
\label{NG01}

\paragraph{Finding NG01.}
Let $M$ be a pure motive of abelian type over $\mathbb Q$, for absolute
Hodge cycles.  There is a canonical decomposition over $\mathbb Q$
\[
                    M=C(M)\oplus N(M),                         \tag{NG01.1}
\]
where $C(M)$ is potentially CM and $N(M)$ has no nonzero potentially-CM
motivic subobject.  Let $G_M$ be the algebraic group of the neutral
Tannakian category generated by $M$, for the Betti fibre functor, and put
$D_M=(G_M^\circ)^{\rm der}$.  Then
\[
 H_B(C(M))=H_B(M)^{D_M},\qquad
 H_B(N(M))=\langle (d-1)v:d\in D_M,\ v\in H_B(M)\rangle.       \tag{NG01.2}
\]
The second expression is interpreted after scalar extension to an
algebraic closure and then descended.  Equivalently it is the sum of the
nontrivial $D_M$-isotypic summands.  These formulas retain the action of
the entire, possibly disconnected group $G_M$.

Every $\ell$-adic realization of (NG01.1) is an actual direct sum of
Galois representations.  When $M=h^j(A)$, its CM part is therefore
available for the Hecke factorization of ZETA\_TWO CM01, and the quotient
of the good-prime local polynomial of $h^j(A)$ by that of $C(M)$ is the
local polynomial of $N(M)$.  This gives an exact factorization of the
incomplete $L$-function.

\paragraph{Proof.}
The category is semisimple neutral Tannakian by Deligne--Milne,
Proposition~6.5 \cite{NGDeligneMilne}.  The fully faithful Hodge realization over an
algebraically closed field (Theorem~6.25) and the descent sequence
(Proposition~6.23) identify $G_M^\circ$ with the Mumford--Tate group of
$H_B(M)$; see also Commelin, Remark~4.14(1) \cite{NGCommelin}.  In particular $G_M^\circ$
is reductive.  Its derived group $D_M$ is connected semisimple and is
characteristic in $G_M^\circ$, hence normal in $G_M$.

Restrict the representation $V=H_B(M)$ to $D_M$.  Complete reducibility
gives the direct sum of its trivial isotypic component $V^{D_M}$ and
the sum $V_{\ne1}$ of all its nontrivial isotypic components.  Normality
of $D_M$ implies that every $g\in G_M$ preserves the trivial isotypic
component and permutes the nontrivial ones.  Thus both summands and the
projection $e:V\to V^{D_M}$ are $G_M$-equivariant.  Tannaka duality
identifies $e$ with an idempotent endomorphism of $M$ over $\mathbb Q$.
This proves the descent of the projector, rather than merely its
existence after an unspecified finite extension.  Its image and kernel
give (NG01.1), and semisimplicity gives (NG01.2).

The identity component of the group acting on $V^{D_M}$ is a quotient
of the torus $G_M^\circ/D_M$.  An abelian motive with torus identity
component is in the category generated by Artin motives and potentially
CM abelian motives (Commelin, \S4.16, following Lemma~3.8 and
Theorem~4.13).  Hence $C(M)$ is potentially CM.  Conversely the connected
group acting on a potentially-CM subobject is a torus.  A morphism from
the connected semisimple group $D_M$ to this torus is trivial.  Such a
subobject is therefore contained in $V^{D_M}$, which proves maximality
and the assertion about $N(M)$.  The same argument proves uniqueness
and functoriality for morphisms: apply it in the common category
generated by source and target, where a morphism respects the trivial
and nontrivial isotypic components.  The formula using spans agrees with
$V_{\ne1}$ because its quotient is $D_M$-trivial and the category of
$D_M$-representations is semisimple.

A realization functor preserves direct sums and idempotents, so the
$\ell$-adic splitting follows.  At a good prime the determinant of
Frobenius on a direct sum is the product of the two determinants.  The
CM factor has the compatible Hecke description already used in
ZETA\_TWO CM01.  Division of the known good-prime polynomial of $h^j(A)$
by that polynomial is therefore a polynomial, independent of the
auxiliary $\ell$, and is the polynomial on $N(M)_\ell$.  This last
argument avoids assuming a separate compatibility theorem for every
possible absolute-Hodge projector.  Multiplication of the local
identities proves the asserted incomplete Euler-product identity.
\hfill$\square$

\paragraph{The precise role of Faltings.}
Faltings, \S5, S\"atze~3--4, printed p.~360 \cite{NGFaltings}, proves semisimplicity of
$H^1_\ell(A)$ and the endomorphism theorem.  Since the Zariski closure
$G_\ell$ acts faithfully on this semisimple representation, its
unipotent radical is trivial; in characteristic zero all its tensor
and exterior representations are semisimple.  Consequently
$H^j_\ell(A)=\bigwedge^jH^1_\ell(A)$ is semisimple for every $j$.
For clarity, the reductivity implication can be checked directly:
on every irreducible constituent the normal connected unipotent
radical has a nonzero fixed subspace, which is stable under $G_\ell$,
and therefore is the whole constituent.  Faithfulness forces that
radical to be trivial.  The finite component group causes no
semisimplicity problem in characteristic zero.

At a fixed $\ell$ there is also a canonical maximal potentially
abelian subrepresentation
\[
 \bigl(H^j_\ell(A)\bigr)^{(G_\ell^\circ)^{\rm der}}.           \tag{NG01.3}
\]
One must not identify (NG01.3) with $C(h^j(A))_\ell$ by citing
Faltings alone.  Deligne gives the inclusion
$G_\ell^\circ\subseteq\operatorname{MT}(A)_{\mathbb Q_\ell}$,
and hence inclusion of the motivic CM part in (NG01.3).
Equality follows, for example, when the Mumford--Tate conjecture is
known for $A$.  Faltings' endomorphism theorem for $H^1$ does not assert
that every projector in a higher tensor power is an absolute-Hodge
projector.  The proof of NG01 uses the motivic group and supplies the
required motivic splitting without making this identification.

\paragraph{Scope of the closure.}
NG01 closes the decomposition gap in ZETA\_TWO ZT06 whenever the pure
factor is of abelian type.  Deligne--Milne, Proposition~6.26 and
Corollary~6.27,
include curves, products with abelian and zero-dimensional varieties,
Fermat hypersurfaces, K3 surfaces, and the stated unirational varieties
of dimension at most three.  The conclusion is about motives for
absolute Hodge cycles.  It is not a general construction of Chow
projectors, and it does not assert such a splitting for an arbitrary
smooth projective variety outside the abelian-type category.

\subsection{NG02: the exact list for every curve}
\label{NG02}

\paragraph{Finding NG02.}
Let $C/\mathbb Q$ be a reduced separated finite-type curve.  Zero-dimensional
components can be added as separate Artin terms, so first assume its
components have dimension one.  Passing to the reduction does not
change the zeta function, so this also treats nonreduced curves.
Let $\widetilde C\to C$ be its
normalization, $\overline C$ the smooth projective compactification of
$\widetilde C$, $D=\overline C\setminus\widetilde C$, $Z$ the finite
singular locus of $C$, and $\widetilde Z$ its inverse image.  All finite
schemes in the following formula are reduced.  Put
\[
 \Pi=H^0(\overline C_{\overline{\mathbb Q}},\mathbb Q),\qquad
 B=H^0(Z)-H^0(\widetilde Z)-H^0(D),\qquad
 J=\operatorname{Pic}^0(\overline C/\mathbb Q).
\]
Thus $\Pi$ is the permutation Artin motive of geometric connected
components and $B$ is a virtual Artin motive.  For a sufficiently
large finite $S$, the incomplete zeta function of a spread-out model is
\[
 \zeta^S_C(s)=
 \frac{L^S(\Pi+B,s)L^S(\Pi,s-1)}{L^S(H^1(J),s)}.             \tag{NG02.1}
\]
Write the $\mathbb Q$-isogeny decomposition
\[
                     J\sim_{\mathbb Q}\prod_i B_i^{m_i},
\]
with $B_i$ pairwise nonisogenous $\mathbb Q$-simple abelian varieties,
and let $I_{\rm nc}$ consist of those that are not geometrically of CM
type.  The complete non-CM list in (NG02.1) is precisely
\[
            \bigl\{\,L^S(B_i,s)\text{ with multiplicity }m_i:
                                i\in I_{\rm nc}\,\bigr\}.     \tag{NG02.2}
\]
Each factor has weight one and its desired RH line is $\Re s=1$.
All corrections for singularities, missing points, and geometric
components are already in the explicitly displayed Artin factors.

\paragraph{Proof.}
Normalization is an isomorphism away from the indicated finite closed
sets.  The scissor relation gives
\[
 [C]=[\overline C]-[D]-[\widetilde Z]+[Z]
                  \quad\text{in }K_0(\operatorname{Var}_{\mathbb Q}).
\]
The relations spread over a sufficiently small open of
$\operatorname{Spec}\mathbb Z$ by the same finite-presentation argument
used in ZETA\_TWO ZT06 (EGA IV$_3$, (8.8.2)).  For a smooth projective
curve the pure cohomology is $\Pi$, $H^1(J)$, and $\Pi(-1)$ in degrees
zero, one, and two.  The geometrically connected case is
Deligne--Milne, Proposition~6.14; for general components the same
identity descends with the permutation action, or equivalently follows
by restriction of scalars from their fields of definition.  No
rational base point is needed: changing a geometric Abel--Jacobi map
by translation does not change its map on $H^1$.  Applying
the trace formula to the spread-out scissor relation proves (NG02.1).

The isogeny decomposition gives the direct sum of the $H^1(B_i)$ with
the stated multiplicities.  The sum of all isotypic factors that
become CM is stable under Galois, so its central idempotent is defined
over $\mathbb Q$.  Equivalently it is the degree-one projector of NG01.
Deligne--Milne, Proposition~6.21, identifies degree-one motives with
abelian varieties up to isogeny; therefore a nonzero CM submotive of
$H^1(B_i)$ would give a nonzero CM isogeny factor of $B_i$.  Simplicity
forces the whole $B_i$ to have potential CM.  Thus no such subobject
exists for $i\in I_{\rm nc}$, and every other factor belongs to CM01.
This proves completeness of (NG02.2).  Faltings' theorem verifies the
corresponding degree-one semisimple Galois decomposition, including
the endomorphism idempotents.  No splitting of the mixed
compact-support cohomology of an open or singular curve has been
assumed; its Grothendieck class is what the scissor calculation uses.
\hfill$\square$

\paragraph{Extension to nonseparated curves.}
The same exact non-CM list applies to any finite-type
$\mathbb Q$-scheme $C$ of dimension at most one, without a
separatedness assumption.  Reduce $C$ first.  For every
one-dimensional irreducible component choose a dense affine open
$U_i$ contained in that component and avoiding all other components.
These opens are also open in $C$ and are pairwise disjoint: take
an affine neighbourhood of its generic point inside the open
complement of the other components.  Put $U=\coprod_iU_i$ and
$Z_0=(C\setminus U)_{\rm red}$.  Every one-dimensional generic point
lies in $U$, so $Z_0$ is zero-dimensional and finite over $\mathbb Q$.
For this last assertion no separatedness is needed: a zero-dimensional
finite-type scheme has finitely many open-and-closed points, each
with an Artinian finite-dimensional coordinate algebra over the base.
Thus $U$ is separated, and the closed-point scissor identity gives,
after enlarging $S$,
\[
               \zeta_C^S(s)=\zeta_U^S(s)L^S(H^0(Z_0),s).
                                                               \tag{NG02.4}
\]
Apply (NG02.1) to $U$: its virtual Artin term $B_U$ is simply
replaced by $B_U+H^0(Z_0)$.  The residual list (NG02.2) and the
order identity (NG02.3) are unchanged with this replacement.
Indeed the added factor is Artin, and the smooth projective
normalizations of the $U_i$ depend only on the function fields of
the original components.  Consequently the Jacobian factors in
the list do not depend on the affine opens chosen.  If $U$ is empty,
the entire zeta function is Artin.  This proves the extension by
separated strata and does not require a proper compactification
of a nonseparated normalization.\hfill$\square$

\paragraph{A precise coefficient refinement.}
For a $\mathbb Q$-simple $B_i$, let $E_i$ be the centre of the division
algebra $\operatorname{End}^0_{\mathbb Q}(B_i)$ and write its dimension
over $E_i$ as $d_i^2$.  Over $\overline{\mathbb Q}_\ell$, Faltings'
endomorphism theorem and the double-centralizer decomposition give
$[E_i:\mathbb Q]$ pairwise nonisomorphic irreducible constituents, each
of rank
\[
                  r_i=\frac{2\dim B_i}{d_i[E_i:\mathbb Q]},
\]
and each of multiplicity $d_i$.  Indeed the rational Betti realization
is a free module over the division algebra, so its dimensions at all
embeddings of $E_i$ are equal.  The commutant becomes the
product of $[E_i:\mathbb Q]$ matrix algebras
$M_{d_i}(\overline{\mathbb Q}_\ell)$; comparison with the commutant of
a semisimple representation proves both assertions.  This identifies
the degree-two (GL$_2$-type) systems and the higher-rank systems in
(NG02.2) without calling every abelian-variety factor a GL$_2$
$L$-function.  If a coefficient refinement or automorphic realization
is known, it may be used; the rational $L(B_i,s)$ is already an exact
named factor and needs no such additional identification.

\paragraph{Noncancellation for the curve product.}
Under \textbf{M}, and assuming meromorphic continuation of the displayed
non-CM product, every point with
$\tfrac12<\Re s<\tfrac32$ and $\Re s\ne1$ satisfies
\[
 \operatorname{ord}_s\zeta^S_C
       =-\operatorname{ord}_s
             \prod_{i\in I_{\rm nc}}L^S(B_i,s)^{m_i}.          \tag{NG02.3}
\]
To check this, the weight-zero Artin factors have their nontrivial
divisors on the boundary $\Re s=1/2$, and their weight-two twists on
the boundary $\Re s=3/2$.  Their possible real endpoint divisors in
this open strip are on $\Re s=1$.  The CM weight-one factors have
their global divisors on $\Re s=1$ by CM01; their omitted local
Hecke factors have weight one and therefore divisors on
$\Re s=1/2$.  Deleted factors of the Artin terms lie on real parts
zero or one.  Thus all other terms are invertible at the indicated
points and the order identity follows directly from (NG02.1).
This version with incomplete functions needs no new assertion about
the complete bad-prime factors of a general non-CM abelian variety.
When individual non-CM factors are known to be holomorphic
in the strip, their positive multiplicities prevent cancellation
among them.  Without holomorphy, the exact converse is for the
aggregate meromorphic product, not for every separate factor.

\subsection{NG03: the exact abelian-variety list}
\label{NG03}

\paragraph{Finding NG03.}
Let $A/\mathbb Q$ have dimension $g$ and choose a polarization over
$\mathbb Q$.  Let $P^r(A)$ denote primitive degree-$r$ cohomology,
$0\le r\le g$, and set
\[
                  C_r=C(P^r(A)),\qquad N_r=N(P^r(A)).
\]
These are absolute-Hodge motives over $\mathbb Q$.  Then
\[
 \zeta_A^S(s)=
 \prod_{r=0}^g\prod_{k=0}^{g-r}
 \left\{L^S(C_r,s-k)L^S(N_r,s-k)\right\}^{(-1)^r}.            \tag{NG03.1}
\]
Here $P^0=C_0=\mathbb Q(0)$ and $N_0=0$.  Consequently the exact
residual list after the CM part of ZT06 is removed is
\[
 \boxed{\ \left\{L^S(N_r,s-k):1\le r\le g,\quad
                    0\le k\le g-r,\quad N_r\ne0\right\}.\ } \tag{NG03.2}
\]
The exponent is $(-1)^r$, the weight is $r+2k$, and the desired line is
$\Re s=(r+1)/2+k$.  The remaining analytic requirements are precisely
continuation, functional equation and divisor RH for these factors,
to the extent not known independently, together with the explicit
finite local factors of ZT04.  Under \textbf{M} the CM factors satisfy
the conclusion of CM01; no assertion that \textbf{M} settles
(NG03.2) is made.

\paragraph{Proof.}
The cohomology algebra is the exterior algebra of $H^1(A)$
(Deligne--Milne, Proposition~6.13).  The symplectic Lefschetz
decomposition is
\[
 h^j(A)=\bigoplus_{\substack{r+2k=j\\0\le r\le g\\0\le k\le g-r}}
                         P^r(A)(-k).                         \tag{NG03.3}
\]
One may construct it directly on the exterior algebra: in a symplectic
basis, wedge multiplication by the polarization and contraction by
the inverse alternating form, with the degree operator, form the
standard $\mathfrak{sl}_2$ triple.  Every irreducible
$\mathfrak{sl}_2$-string beginning in degree $r\le g$ has the degrees
$r,r+2,\ldots,2g-r$.  Its lowest-degree vectors are exactly the kernel
of contraction, namely $P^r$; successive multiplication by the
polarization supplies the indicated Tate twists.  Complete
reducibility gives (NG03.3).  The operators are Hodge morphisms and
are absolute Hodge by Deligne's theorem, so the decomposition exists
in the chosen motivic category over $\mathbb Q$.

Apply NG01 to each primitive object.  Tate twists do not change
whether the connected derived group acts trivially, hence
$C(P^r(-k))=C_r(-k)$ and $N(P^r(-k))=N_r(-k)$.  Frobenius determinants
multiply on these direct sums.  Degree $j=r+2k$ has parity $r$, and
$L(M(-k),s)=L(M,s-k)$.  The trace-formula Euler product now gives
(NG03.1).  The list (NG03.2) and its weight and signs follow.  As in
ZT06, this is an exact factorization and a sufficient analytic gate;
it does not claim a converse that would rule out meromorphic
cancellation between distinct residual factors.\hfill$\square$

\paragraph{What is, and is not, computed uniformly.}
For a specified Mumford--Tate representation $V=H^1_B(A)$, (NG03.2)
is finite linear algebra.  If $X_1,\ldots,X_t$ span
$\operatorname{Lie}(\operatorname{MT}(A)^{\rm der})$, then
\[
             H_B(C_r)=\bigcap_{a=1}^t
                  \ker\bigl(d\rho_r(X_a):P^r_B\to P^r_B\bigr).
                                                                    \tag{NG03.4}
\]
The derived group is connected, so Lie-algebra invariants and group
invariants agree.  The complementary isotypic subspace then gives
$N_r$, and local determinants give its Euler factors.  No general
effective algorithm recovering the full Mumford--Tate representation
of an arbitrary abelian variety from equations is proved by
Faltings or Deligne.  Thus (NG03.2)--(NG03.4) are an exact universal
formula with a specified input, not a claimed algorithm solving that
separate identification problem.  If only $A\sim A_{\rm CM}\times
A_{\rm nc}$ is known, the correct formula is
\[
 C(h^j(A))=
 \bigoplus_{a+b=j}\bigwedge^a h^1(A_{\rm CM})\otimes
                       C\!\left(\bigwedge^b h^1(A_{\rm nc})\right).
                                                                    \tag{NG03.5}
\]
Indeed the derived group acts trivially on the first factor and
surjects onto the derived group acting on the second.  This formula
does not discard the CM summands created inside exterior powers of
the non-CM part.

\subsection{NG04: the full symplectic case}
\label{NG04}

\paragraph{Finding NG04.}
Suppose $\operatorname{MT}(A)=\operatorname{GSp}_{2g}$ in its standard
representation.  Then $C(h^{2i}(A))=\mathbb Q(-i)$,
$C(h^{2i+1}(A))=0$, and $N_r=P^r$ for $1\le r\le g$.
In particular
\[
 \zeta_A^S(s)=\prod_{i=0}^g\zeta^S(s-i)
       \prod_{r=1}^g\prod_{k=0}^{g-r}L^S(P^r(A),s-k)^{(-1)^r},
 \qquad \operatorname{rank}P^r={2g\choose r}-{2g\choose r-2}. \tag{NG04.1}
\]
The non-CM functions are the primitive exterior-power Hasse--Weil
$L$-functions and their indicated shifts.  For a full-symplectic
abelian surface this means $L(H^1(A),s)$,
$L(H^1(A),s-1)$, and the rank-five primitive exterior-square function
$L(P^2(A),s)$; the lines are $1$, $2$, and $3/2$, respectively.

\paragraph{Proof.}
The derived group is $\operatorname{Sp}_{2g}$.  We give a direct check
of its invariants in the exterior algebra, so no assumption about
unknown algebraic cycles is involved.  Choose a symplectic basis
$e_1,f_1,\ldots,e_g,f_g$.  The diagonal torus sends $e_i$ to $t_i e_i$
and $f_i$ to $t_i^{-1}f_i$.  A zero-weight exterior monomial must
contain both $e_i,f_i$ or neither, for every $i$; hence no odd-degree
invariant exists.  In degree $2i$, the zero-weight monomials are the
products of $i$ distinct pairs $e_a\wedge f_a$.  The subgroup
permuting the symplectic pairs acts transitively on these monomials,
with positive sign because each pair has degree two.  Thus any
$\operatorname{Sp}_{2g}$-invariant is a scalar multiple of their sum.
That sum is $1/i!$ times the $i$-th power of the polarization bivector,
which is indeed invariant.  The invariant space therefore has
dimension exactly one.  NG01 identifies it with $\mathbb Q(-i)$,
using the polarization defined over $\mathbb Q$.

For $r\le g$, wedge multiplication by the polarization is injective
from degree $r-2$ to degree $r$.  Its image and the primitive kernel
form a direct sum by the symplectic Lefschetz decomposition in NG03.
This gives the displayed rank.  In positive even primitive degree
the sole invariant, the power of the polarization, is in the image
of wedge multiplication and hence has zero intersection with the
primitive kernel.  In odd degree there is no invariant at all.
Consequently $C(P^r)=0$ for every $r\ge1$.  Substitute in NG03.
\hfill$\square$

\subsection{NG05: every power of a non-CM elliptic curve}
\label{NG05}

\paragraph{Finding NG05.}
Let $E/\mathbb Q$ be non-CM and $m\ge1$.  Adopt the convention that
$\binom m t=0$ outside $0\le t\le m$.  For $0\le b\le a\le m$, put
\[
 d_m(a,b)=\binom ma\binom mb-
                        \binom m{a+1}\binom m{b-1}.           \tag{NG05.1}
\]
The complete decomposition is
\[
 h^j(E^m)=\bigoplus_{\substack{0\le b\le a\le m\\a+b=j}}
       \left(\operatorname{Sym}^{a-b}h^1(E)(-b)\right)^{\oplus d_m(a,b)}.
                                                                    \tag{NG05.2}
\]
The CM terms are exactly $a=b$, and all other terms are non-CM.
Consequently
\[
 \zeta_{E^m}^S(s)=
 \prod_{0\le b\le a\le m}
 L^S(\operatorname{Sym}^{a-b}E,s-b)^{(-1)^{a+b}d_m(a,b)},       \tag{NG05.3}
\]
where $L(\operatorname{Sym}^0E,s)=\zeta(s)$.  The distinct non-CM
analytic inputs, before their specified shifts, are exactly
$L(\operatorname{Sym}^nE,s)$, $1\le n\le m$.  Their motivic RH lines
are $\Re s=(n+1)/2$.  Their analytic continuation is already known:
for $n\ge2$ see Newton--Thorne, \emph{Symmetric power functoriality
for holomorphic modular forms, II}, Theorem~A and Corollary~B
(printed pp.~117--118) \cite{NGNewtonThorne}; the $n=1$ case is the modularity theorem
used in ZETA\_TWO ZT07.  Their RH is an additional analytic assertion.

\paragraph{Proof.}
Write $V=h^1(E)$ and $W=\mathbb Q^m$ with trivial action.
Then $h^j(E^m)=\bigwedge^j(V\otimes W)$.  The Mumford--Tate group
of a non-CM elliptic curve is $\operatorname{GL}_2$ (for example,
Moonen, Example~2.1.13 \cite{NGMoonen}); its derived group is $\operatorname{SL}_2$.
We compute the decomposition directly.  On a diagonal matrix with
eigenvalues $x,y$, the exterior-character generating series is
\[
                      (1+tx)^m(1+ty)^m.
\]
The coefficient of $t^{a+b}x^ay^b$ is
$w(a,b)=\binom ma\binom mb$.  The irreducible
$\operatorname{GL}_2$-representation of highest weight $(a,b)$,
$a\ge b$, has character
$x^ay^b+x^{a-1}y^{b+1}+\cdots+x^by^a$ and is
$\operatorname{Sym}^{a-b}V\otimes(\det V)^b$.
For fixed total degree, subtraction of the adjacent weight
multiplicities $w(a,b)-w(a+1,b-1)$ therefore isolates its highest-weight
multiplicity.  This is (NG05.1).  The exterior representation is
semisimple, so this character calculation gives its direct-sum
decomposition.  Equivalently, it is the two-row case of the exterior
Cauchy decomposition.  The corresponding Hodge-equivariant
projectors are absolute Hodge and descend as morphisms in the
$\mathbb Q$-category generated by $V$: its full motivic group is a
subgroup of $\operatorname{GL}(V)$ whose identity component is already
$\operatorname{GL}(V)$, so it equals that group.  Thus the identity is motivic.
Since $\det V=\mathbb Q(-1)$, it is exactly (NG05.2).

For $a>b$, $\operatorname{Sym}^{a-b}V$ is a nontrivial irreducible
$\operatorname{SL}_2$-representation, hence has no derived-group
invariants.  For $a=b$ it is the trivial representation, followed by
a Tate twist.  NG01 proves the asserted complete CM/non-CM
classification.  Taking determinants in degree $j=a+b$ gives
(NG05.3).\hfill$\square$

In particular the first example beyond a curve is
\[
 \zeta_{E^2}^S(s)=
 \frac{\zeta^S(s)\,\zeta^S(s-1)^3\,\zeta^S(s-2)
                        L^S(\operatorname{Sym}^2E,s)}
      {L^S(E,s)^2 L^S(E,s-1)^2}.                              \tag{NG05.4}
\]
Thus the non-CM gate for $E^2$ contains a symmetric-square
$L$-function, even though the only simple abelian isogeny factor is
$E$.  Conversely its degree-two cohomology has three Tate summands,
even though $h^1(E^2)$ has no CM summand.  This refutes the procedure
of removing only the exterior algebra of the CM isogeny factors and
calling everything left non-CM.

\subsection{NG06: non-Tate CM terms can occur in higher degree}
\label{NG06}

\paragraph{Finding NG06.}
Let $A/\mathbb Q$ have dimension $g$, and suppose that over a number
field $F$ it has an action by an imaginary quadratic field $K$.
Suppose that the $K$-action on $H^1(A)$ has signature $(p,q)$, where
$p+q=g$.  Over $F$, the rational Hodge subspace
\[
                         W=\bigwedge_K^g H^1(A)               \tag{NG06.1}
\]
has rank two over $\mathbb Q$, occurs as a summand of $H^g(A)$,
and has Hodge types $(p,q)$ and $(q,p)$.  It is a CM absolute-Hodge
motive.  If $p\ne q$, it is not Artin--Tate.  Its Galois-stable sum of
conjugates in $h^g(A)_{\overline{\mathbb Q}}$ descends to a CM summand
over $\mathbb Q$, and is included by NG01.

\paragraph{Proof.}
Over a splitting field for $K$, write
$H^1=V_\sigma\oplus V_{\bar\sigma}$, with each summand of dimension
$g$.  The exterior algebra decomposes by its bidegree in these two
spaces.  The sum
\[
                      \bigwedge^g V_\sigma\ \oplus\
                      \bigwedge^g V_{\bar\sigma}
\]
is invariant under the coefficient Galois group and defines the
rational direct summand (NG06.1).  Taking determinants adds the Hodge
indices, so its two types are precisely $(p,q)$ and $(q,p)$.
The projector onto these two extreme exterior degrees, along all
the mixed degrees, is invariant under $G_F$ because the $K$-action
is defined over $F$.  Its absolute-Hodge realization thus descends
over $F$.
The summand is a free rank-one $K$-module.  Its Mumford--Tate group
commutes with $K$ and is consequently contained in
$\operatorname{Res}_{K/\mathbb Q}\mathbb G_m$.  It is therefore a
torus.  Deligne's theorem makes its Hodge projector an absolute-Hodge
projector, and the CM characterization used in NG01 proves the CM
assertion.  An Artin--Tate pure Hodge structure has only diagonal
types, so $p\ne q$ excludes that possibility.

The projector is defined over a finite extension by Deligne,
Proposition~2.9 \cite{NGDeligne82}.  There are only finitely many conjugate subobjects.
Their sum is Galois stable, remains CM because the category is closed
under direct sums and quotients, and has a projector in the semisimple
category.  Conjugation preserves the CM category because it preserves
Artin motives and the endomorphism criterion for CM abelian varieties.
Choose a finite Galois extension defining the sum and a projector
onto it, and average the projector over its finite Galois group.
Every conjugate projector has image in the same subobject and is the
identity on it.  Their average consequently has the same image and
is still idempotent; it is Galois invariant and descends over
$\mathbb Q$.  This proves the final assertion without assuming the
original $K$-action is defined over $\mathbb Q$.\hfill$\square$

For example, signature $(2,1)$ in dimension three yields a weight-three
CM term of types $(2,1)$ and $(1,2)$.  The statement is a construction
for every abelian variety with the specified action and signature;
it does not claim that an arbitrary abelian threefold has that action.
Together with NG05, it explains why higher CM terms must be selected
degree by degree and cannot be read solely from a list of CM
degree-one factors.

\subsection{NG07: an explicit non-CM Jacobian with non-Tate CM cohomology}
\label{NG07}

\paragraph{Finding NG07.}
Let $C/\mathbb Q$ be the smooth projective curve with affine equation
\[
                            y^3=x^5-x-1,
\]
and let $A=\operatorname{Jac}(C)$.  Then $A$ is an absolutely simple
non-CM abelian fourfold and $h^1(A)$ has no nonzero CM summand.
Nevertheless $h^4(A)$ contains a rank-two CM summand $W$ defined over
$\mathbb Q$, with Hodge types $(3,1)$ and $(1,3)$.  In particular
$W$ is not Artin--Tate.  It occurs in the primitive part $P^4(A)$
for the canonical Jacobian polarization.

\paragraph{Proof.}
Zarhin, Theorem~1.1 and Examples~1.3(1) of
\emph{The endomorphism rings of jacobians of cyclic covers of the
projective line} \cite{NGZarhin}, apply to the exact family
$y^p=x^n-x-1$ with $n\ge5$ and $p$ odd.  For $n=5,p=3$ they give
\[
 \operatorname{End}_{\overline{\mathbb Q}}(A)=\mathbb Z[\zeta_3],
 \qquad \dim A=(3-1)(5-1)/2=4.
\]
In particular $A$ is absolutely simple.  Its geometric endomorphism
algebra has dimension two over $\mathbb Q$, so cannot contain a
commutative semisimple subalgebra of dimension $2\dim A=8$; thus
$A$ is non-CM.  NG02, or the degree-one equivalence in
Deligne--Milne 6.21, shows that $h^1(A)$ has no CM subobject.
For separate input checks, the discriminant of $x^5-x-1$ was computed as
$2869$, its reduction modulo $3$ is irreducible, and its reduction
modulo $2$ is
$(x^2+x+1)(x^3+x^2+1)$.  These computations are recorded separately;
the Galois-group conclusion used here is Zarhin's stated example.

Put $K=\mathbb Q(\zeta_3)$ and let
$\delta(x,y)=(x,\zeta_3y)$.  A basis of regular differentials is
\[
 \frac{dx}{y},\quad \frac{dx}{y^2},\quad
                    \frac{x\,dx}{y^2},\quad
                    \frac{x^2\,dx}{y^2}.                    \tag{NG07.1}
\]
For completeness, at a finite branch point $y$ is a local parameter
and $dx$ has order two, so all four forms are regular.  At the unique
point at infinity, $x,y,dx$ have orders $-3,-5,-4$, respectively.
The forms in (NG07.1) therefore have orders $1,6,3,0$, respectively,
and are regular there.  They are linearly independent: the first has
$\delta$-eigenvalue $\zeta_3^2$, while the remaining three have
eigenvalue $\zeta_3$ and are independent powers of $x$ times the same
nonzero differential.  The genus is four, so they form a basis.
This proves that the $K$-signature is $(3,1)$, up to exchanging the
two embeddings.

NG06 now supplies
$W=\bigwedge_K^4 H^1(A)$, of types $(3,1)$ and $(1,3)$.
Here $K$ is the entire geometric endomorphism algebra, so Galois acts
on it by field automorphisms, either fixing or exchanging its two
embeddings.  The sum of the two extreme exterior eigenspaces that
defines $W$ and its complementary mixed-degree subspace are
therefore stable under the full Galois group.  Their absolute-Hodge
projector is defined over $\mathbb Q$, so $W$ itself, with rank two,
descends; it is not necessary to enlarge it to a sum of conjugates.

Finally $\delta$ preserves the canonical Jacobian polarization.
Consequently the polarization pairs the two $K$-eigenspaces in
$H^1$ with one another and has no same-eigenspace component.  The
contraction operator therefore kills both
$\bigwedge^4 V_\sigma$ and $\bigwedge^4 V_{\bar\sigma}$.
This puts $W$ in $P^4(A)$.  Its non-diagonal Hodge types prove that
it is not Artin--Tate.\hfill$\square$

\paragraph{Placement in Grothendieck's frameworks.}
Over $U=\operatorname{Spec}\mathbb Z[1/S\ell]$, for separated models
the formulas are
identities in $K_0(D_c^b(U,\mathbb Q_\ell))$ for the realizations of
$Rf_!\mathbb Q_\ell$, with actual pure direct-sum projectors in the
abelian-type cases and a virtual scissor identity for open or singular
curves.  The nonseparated-curve extension (NG02.4) uses the finite sum
of the compact-support classes of its separated strata.
The support and trace framework is SGA~4 XVII,
\S\S5.1--5.3, and SGA~5 XV, \S3, no.~2, equation~(1), as located and
used in ZETA\_TWO ZT04--ZT06.  Exterior powers, determinants and
their passage to Grothendieck groups are the $\lambda$-ring mechanism
of SGA~6 V, \S2 and \S7, (7.1.1), and ZETA\_TWO ZT08.
The relevant documented archive locations are cote~60, scans~3 and~5
for determinant/exterior-power algebra, scan~7 for reciprocity, and
scan~9 for the explicit Weil and dual/Tate hypotheses.  Cote~34,
scans~64 and~74, remains the finite-field Riemann--Weil context
recorded in the predecessor.  These are carried archival locators,
not new readings of unavailable scans, and none is attributed the
later Faltings--Deligne splitting theorem.  In the runs, NG01--NG03
close the specific abelian-type decomposition gap noted in
INTAKE\_ZETA\_TWO, \S1, row ZT06 and \S5.1; NG05 refines the first
non-CM example ZT07; NG01 and NG06--NG07 extend the usable scope of CM01.
